\documentclass[a4paper,fontset=none]{ctexart}
\usepackage[a4paper,left=17mm,right=17mm,top=14mm,bottom=15mm,footskip=8mm]{geometry}
\usepackage{amsmath,amssymb,mathtools}
\usepackage{booktabs,array,tabularx}
\usepackage{xcolor,fancyhdr,fontspec,microtype}

\setmainfont{Times New Roman}
\ifdefined\DGPFontPath
  \setCJKmainfont[
    Path=\DGPFontPath,
    Extension=.ttf,
    UprightFont=*-Regular,
    BoldFont=*-Bold
  ]{NotoSerifSC}
\else
  \setCJKmainfont[AutoFakeBold=2.5]{STSong}
\fi

\definecolor{ink}{HTML}{202328}
\definecolor{muted}{HTML}{68717A}
\definecolor{rule}{HTML}{8A929A}
\definecolor{soft}{HTML}{F2F3F4}
\color{ink}

\setlength{\parindent}{0pt}
\setlength{\parskip}{1.5pt}
\setlength{\abovedisplayskip}{4pt}
\setlength{\belowdisplayskip}{4pt}
\setlength{\jot}{1pt}
\renewcommand{\arraystretch}{1.17}

\pagestyle{fancy}
\fancyhf{}
\renewcommand{\headrulewidth}{0pt}
\fancyfoot[L]{\fontsize{8}{9}\selectfont\color{muted} AXID-DC v10}
\fancyfoot[R]{\fontsize{8}{9}\selectfont\color{muted}\thepage/4}

\newcommand{\pagetitle}[1]{%
  \begin{center}
    {\fontsize{19}{22}\selectfont\bfseries Empirically Calibrated DGP}\par
  \end{center}
  \vspace{4pt}
  \color{rule}\hrule height .55pt\color{ink}
  \vspace{5pt}
  \noindent\colorbox{soft}{%
    \parbox{\dimexpr\linewidth-2\fboxsep\relax}{%
      \fontsize{13}{15}\selectfont\bfseries #1}}
  \vspace{5pt}}
\newcommand{\group}[1]{%
  \vspace{4pt}{\fontsize{11}{13}\selectfont\bfseries #1}\par
  \vspace{1pt}\color{rule}\hrule height .4pt\color{ink}\vspace{3pt}}
\newcommand{\smallnote}[1]{%
  {\fontsize{8.5}{10}\selectfont\color{muted}#1}}
\newcommand{\TN}{\operatorname{TN}}
\newcommand{\NB}{\operatorname{NB}}
\newcommand{\Beta}{\operatorname{Beta}}
\newcommand{\Gum}{\operatorname{Gumbel}}

\begin{document}
\fontsize{9.4}{11.5}\selectfont

% ============================================================================
% Page 1
% ============================================================================
\pagetitle{第一块：真实-合成特征映射}

\[
\boxed{(u_i,r_i,Hb_i,c_{p,t},c_{e,t},I_t,N_t^-,N_t^+)}
\qquad\text{八类合成生成量}
\]

\group{患者特征}
\begin{center}
\fontsize{8.9}{10.7}\selectfont
\begin{tabularx}{.99\textwidth}{
  >{\raggedright\arraybackslash}p{27mm}
  >{\raggedright\arraybackslash}p{105mm}
  X}
\toprule
合成量 & 真实字段与确定映射 & 模型输入 / 作用\\
\midrule
$u_i$\quad\textbf{申请规模} &
$s_i=\operatorname{coalesce}(\text{申请用血量},0)+
\operatorname{coalesce}(\text{申请备血量},0)$；\newline
$u_i=g(s_i)$，其中
$0/1\mapsto1$，$2$--$6\mapsto s_i$，$7/8\mapsto8$，$\ge9\mapsto10$ &
$u_i/10$\newline 合并稀疏尾部\\

$r_i$\quad\textbf{临床风险} &
$E=\mathbf1\{\text{急救/急诊}\}$，
$B=\mathbf1\{\text{出血/创伤}\}$，\newline
$C=\mathbf1\{\text{白血病/肿瘤/肝病}\}$，
$V=\mathbf1\{\text{年龄}\ge65\}$；
$r_i=(6E+4B+2C+V)/13$ &
输入 $r_i$\newline 仅用决策前字段\\

$Hb_i$\quad\textbf{血红蛋白} &
缺失置训练期中位数 122；非缺失值截断至 $[42,170]$ g/L &
$(Hb_i-42)/128$\\
\bottomrule
\end{tabularx}
\end{center}

\group{日级特征与两个患者数}
\begin{center}
\fontsize{8.9}{10.7}\selectfont
\begin{tabularx}{.99\textwidth}{
  >{\raggedright\arraybackslash}p{27mm}
  >{\raggedright\arraybackslash}p{105mm}
  X}
\toprule
合成量 & 真实字段与确定映射 & 模型输入 / 信息边界\\
\midrule
$c_{p,t}$\quad\textbf{计划压力} &
$n_{p,t}=\#\{\mathcal A_t^-\text{ 中手术备血或病房用血患者}\}$；
$c_{p,t}=\min\{1,\log(1+n_{p,t})/\log89\}$ &
输入 $c_{p,t}$\\

$c_{e,t}$\quad\textbf{急诊压力} &
$n_{e,t}=\#\{\mathcal A_t^-\text{ 中急救用血患者}\}$；
$c_{e,t}=\min\{1,\log(1+n_{e,t})/\log41\}$ &
输入 $c_{e,t}$\\

$I_t$\quad\textbf{期初库存} &
O 型红细胞期初可用库存 &
输入 $I_t/520$\\

$N_t^-$\quad\textbf{可见患者数} &
$N_t^-=|\mathcal A_t^-|$：决策截点前有效申请数 &
输入 $N_t^-/138$\\

$N_t^+$\quad\textbf{新增患者数} &
$N_t^+=|\mathcal A_t^+|$：决策截点后当日新增申请数 &
仅用于需求生成\\
\bottomrule
\end{tabularx}
\end{center}

\group{DAT 的实际输入}
\textbf{患者三维 token}
\[
\boxed{
x_{it}^{\mathrm{pat}}
=\left(\frac{u_i}{10},\,r_i,\,\frac{Hb_i-42}{128}\right)},
\qquad i\in\mathcal A_t^- .
\]
\textbf{日级七维汇总}
\[
\boxed{
X_t=\left(
c_{p,t},c_{e,t},\frac{I_t}{520},\frac{N_t^-}{138},
\frac{\sum_{i\in\mathcal A_t^-}u_i}{1380},
\frac{\sum_{i\in\mathcal A_t^-}r_i}{138},
\frac{\sum_{i\in\mathcal A_t^-}(Hb_i-42)/128}{138}
\right)} .
\]

\newpage

% ============================================================================
% Page 2
% ============================================================================
\pagetitle{第二块：经验校准特征生成}

\begingroup
\fontsize{8.8}{10.4}\selectfont
\renewcommand{\arraystretch}{1.06}
\setlength{\abovedisplayskip}{1.5pt}
\setlength{\belowdisplayskip}{1.5pt}
\setlength{\abovedisplayshortskip}{1pt}
\setlength{\belowdisplayshortskip}{1pt}
\renewcommand{\group}[1]{%
  \vspace{2pt}{\fontsize{10.8}{12.5}\selectfont\bfseries #1}\par
  \vspace{.5pt}\color{rule}\hrule height .4pt\color{ink}\vspace{1.5pt}}

\smallnote{校准期：2013-01-01 至 2019-12-28；2020 年完全留作真实数据测试。合成日期独立抽样。}

\group{日级边际分布}
\begin{center}
\fontsize{9.1}{11}\selectfont
\begin{tabularx}{.96\textwidth}{
  >{\bfseries}p{22mm}
  X}
\toprule
变量 & 冻结边际分布\\
\midrule
$N_t^-$ & $\operatorname{clip}_{[1,138]}\{\NB(6.7049,.105305)\}$\\
$N_t^+$ & $\operatorname{clip}_{[0,73]}\{\NB(4.6768,.142891)\}$\\
$c_{p,t}$ & $\Beta(6.60,1.41)$\\
$c_{e,t}$ & $\Beta(2.95,2.05)$\\
$I_t$ & $.472753\,\TN_{[0,520]}(246.704,37.565^2)
+.527247\,\TN_{[0,520]}(314.756,38.244^2)$\\
\bottomrule
\end{tabularx}
\end{center}

\group{日级经验相关}
\begin{center}
\smallnote{Gaussian copula 潜在相关矩阵；行列顺序均为 $(N^-,N^+,I,c_p,c_e)$。}
\[
\begin{array}{c|rrrrr}
 &N^-&N^+&I&c_p&c_e\\ \hline
N^-&1& .220&-.194& .891& .520\\
N^+&.220&1&-.112& .147& .268\\
I&-.194&-.112&1&-.048&-.439\\
c_p&.891& .147&-.048&1& .240\\
c_e&.520& .268&-.439& .240&1
\end{array}
\]
\end{center}

\group{患者离散边际}
\[
\begin{array}{c|rrrrrrrr}
u&1&2&3&4&5&6&8&10\\ \hline
\Pr(u)&.001219&.220730&.082619&.631536&.003752&.026695&.023018&.010431
\end{array}
\]
\[
\begin{array}{c|rrrrrrr}
k&0&1&2&3&4&5&6\\ \hline
\Pr(r=k/13)&.295695&.191127&.149705&.127568&.012982&.008630&.065153
\end{array}
\]
\[
\begin{array}{c|rrrrrrr}
k&7&8&9&10&11&12&13\\ \hline
\Pr(r=k/13)&.052734&.057030&.018291&.010168&.008217&.002007&.000694
\end{array}
\]

\group{$Hb\mid r$ 的条件混合分布}
令 $L:r<.25$、$M:.25\le r<.55$、$H:r\ge.55$；各成分截断在 $[42,170]$：
\[
\begin{array}{c|c@{\quad}c|c@{\quad}c}
\text{风险层}&w_1&\TN(\mu_1,\sigma_1^2)&w_2&\TN(\mu_2,\sigma_2^2)\\ \hline
L&.201&\TN(79.8,21.0^2)&.799&\TN(135.6,16.0^2)\\
M&.764&\TN(59.8,9.4^2)&.236&\TN(110.2,26.0^2)\\
H&.831&\TN(57.8,8.4^2)&.169&\TN(101.1,27.2^2)\\
\end{array}
\]
\textbf{生成顺序：}
先抽 $r_i$，再抽 $Hb_i^\star\sim Hb\mid r_i$ 并截断至 $[42,170]$；
最后抽 $M_i\sim\operatorname{Bernoulli}(.01506)$，若 $M_i=1$ 则置 $Hb_i=122$，否则 $Hb_i=Hb_i^\star$。

\textbf{患者联合生成：}
$u_i\perp(r_i,Hb_i)$；给定 $N_t^-,N_t^+$ 后，$\mathcal A_t^-$ 与 $\mathcal A_t^+$ 按同一患者画像分布独立抽取。
\endgroup

\newpage

% ============================================================================
% Page 3
% ============================================================================
\pagetitle{第三块：MDCEV 模型}

\group{$\gamma_i$ 与 $\phi_{it}$ 的参数映射}
对当日全部患者 $\mathcal A_t=\mathcal A_t^-\cup\mathcal A_t^+$，
\[
\ell_i=\frac{\log(1+u_i)}{\log11},
\qquad
d_i=\left[\frac{70+10r_i-Hb_i}{10}\right]_+ .
\]
\begin{center}
\fontsize{8.6}{10.2}\selectfont
\begin{tabularx}{.96\textwidth}{
  >{\bfseries}p{14mm}
  X}
$\ell_i$ & 对申请规模作对数变换，平滑右尾并映射到统一尺度。\\
$d_i$ & 风险调整后的 Hb 缺口；$70+10r_i\in[70,80]$ g/L 参考 AABB 讨论的限制性输血阈值范围。\\
\end{tabularx}
\end{center}
\[
\boxed{\gamma_i=.43(.5+1.5\ell_i)},
\qquad
\boxed{\log\phi_{it}=-.8664+.15\ell_i+1.20r_i+.45d_i
+.40c_{p,t}\ell_i+.60c_{e,t}r_i+\eta_t+\epsilon_{it}}.
\]
\[
\eta_t\sim\mathcal N(0,.08^2),
\qquad
\epsilon_{it}\sim\Gum(-.12\gamma_E,.12),
\qquad
\mathbb E[\epsilon_{it}]=0.
\]

\group{系数设置}
\begin{center}
\fontsize{8.8}{10.4}\selectfont
\begin{tabularx}{.98\textwidth}{
  >{\bfseries}p{24mm}
  >{\centering\arraybackslash}p{18mm}
  X}
\toprule
项 & 系数 / 尺度 & 设置含义\\
\midrule
$\gamma$ 尺度 & $.43$ & 控制饱和曲率与数量分配。\\
$\log\phi$ 截距 & $-.8664$ & 与共同库存锚点 $G=280$ 使用一致效用尺度。\\
$\ell_i$ & $.15$ & 申请量已同时进入发血上限和 $\gamma_i$。\\
$r_i$ & $1.20$ & 临床风险的主要患者效应。\\
$d_i$ & $.45$ & 每低于风险调整阈值 10 g/L，缺口增加 1。\\
$c_p\ell_i$ & $.40$ & 计划压力与申请规模交叉项。\\
$c_er_i$ & $.60$ & 急诊压力与临床风险交叉项。\\
$\eta_t,\epsilon_{it}$ & $.08,.12$ & 日级共同扰动与患者个体扰动。\\
\bottomrule
\end{tabularx}
\end{center}

\group{联合优化与需求}
\[
\boxed{
\begin{aligned}
\max_{R,\mathbf Y}\quad
&280\log(R/280)
+\sum_{i\in\mathcal A_t}\gamma_i\phi_{it}
\log\left(1+\frac{Y_i}{\gamma_i}\right)\\
\mathrm{s.t.}\quad
&R+\sum_{i\in\mathcal A_t}Y_i=I_t,
\qquad R>0,
\qquad 0\le Y_i\le u_i.
\end{aligned}}
\]
\[
\text{KKT：}\qquad
Y_i(\lambda_t)=
\min\left\{u_i,\gamma_i\left(\frac{\phi_{it}}{\lambda_t}-1\right)_+\right\},
\]
\[
\frac{280}{\lambda_t}+\sum_iY_i(\lambda_t)=I_t,
\qquad
D_t=\sum_{i\in\mathcal A_t}Y_{it}.
\]

\group{两类交互}
\smallnote{库存与患者集合共同形成影子价格，再由影子价格联合决定个体发血量。}
\[
I_t,\ \{u_j,r_j,Hb_j\}_{j\in\mathcal A_t}
\longrightarrow\lambda_t
\longrightarrow Y_{it}
\qquad\text{\bfseries 患者-患者交互},
\]
\[
(c_{p,t}\ell_i,\,c_{e,t}r_i)
\longrightarrow\phi_{it}
\longrightarrow Y_{it}
\qquad\text{\bfseries 日级状态-患者交互}.
\]

\newpage

% ============================================================================
% Page 4 -- Appendix
% ============================================================================
\pagetitle{附录：经验校准口径}

\smallnote{所有经验量只使用 2013-01-01 至 2019-12-28 的训练期数据；2020 年不参与分布、尺度或目标库存的确定。}

\group{边际分布：为什么这样选，数值如何得到}
\begin{center}
\fontsize{8.1}{9.7}\selectfont
\begin{tabularx}{.99\textwidth}{
  >{\bfseries\raggedright\arraybackslash}p{23mm}
  >{\raggedright\arraybackslash}p{48mm}
  >{\raggedright\arraybackslash}X}
\toprule
对象 & 分布形状的选择 & 参数来源\\
\midrule
$N^-,N^+$ &
日计数方差显著大于均值，采用负二项保留过度离散。 &
矩匹配 $p=m/v$、$a=m^2/(v-m)$：
$N^-:(m,v)=(56.9659,540.9597)$；
$N^+:(28.0529,196.3235)$。
上界 138、73 为训练期最大值。\\

$c_p,c_e$ &
人数经 0.99 分位数对数压缩后位于 $[0,1]$，采用 Beta 描述有界偏态。 &
$c_p$ 使用 $Q_{.99}(n_p)=88$，经验 $(m,v)=(.82449,.01607)$；
$c_e$ 使用 $Q_{.99}(n_e)=40$，经验 $(m,v)=(.58987,.04033)$。
以 Beta 矩公式得到 $(6.60,1.41)$、$(2.95,2.05)$。\\

$I$ &
单一正态不足以保留库存的双状态形状，采用两成分高斯混合；生成时限制在可行支持。 &
2,553 个训练日的两成分 Gaussian-mixture 拟合：
$w=(.472753,.527247)$，
$\mu=(246.704,314.756)$，
$\sigma=(37.565,38.244)$。
训练期最大值 519 向上取整为支持上界 520。\\

$u,r$ &
两者均为有限离散量，不再强加连续分布。 &
$u$ 对总申请量合并稀疏尾部后计算频率；
$r=(6E+4B+2C+V)/13$，直接使用 $k/13$ 的训练期经验频率。\\

$Hb\mid r$ &
风险层内仍有明显异质性，采用两成分混合；生成时将成分截断至经验支持。 &
53,305 条训练申请中，52,502 条 Hb 非缺失；按 $L/M/H$ 分层后，
分别拟合两成分 Gaussian-mixture，得到主文 $w,\mu,\sigma$。
$Q_{.01}=42$、$Q_{.99}=170$、中位数 122；缺失率 $803/53305=.01506$。\\

日级依赖 &
边际分布分别校准，Gaussian copula 只负责保留跨变量依赖。 &
将训练期秩转换为潜在高斯相关，再作
$\Sigma=.95\widehat\Sigma_{\rm rank}+.05I_5$ 收缩；主文报告冻结矩阵。\\
\bottomrule
\end{tabularx}
\end{center}

\group{共同库存锚点与 MDCEV 参数}
\[
G=10\times\operatorname{round}
\left\{\frac{\operatorname{median}(I_{\rm train})}{10}\right\}
=10\times\operatorname{round}(283/10)=280.
\]
\begin{center}
\fontsize{8.5}{10.2}\selectfont
\begin{tabularx}{.98\textwidth}{
  >{\bfseries\raggedright\arraybackslash}p{33mm}
  >{\raggedright\arraybackslash}X}
共同库存水平 & $G=280$ 同时进入医院库存保留效用与血液中心期末库存目标；它是历史库存中位数的业务化取整，不宣称为理论最优安全库存。\\
效用尺度调整 & 固定 $\kappa=1$。原 $G=256$ 下 $(s_\gamma,\beta_0)=(.43,-.956)$ 匹配训练期需求矩 $(20.6227,10.2746,34.0)$；改为 $G=280$ 后取
$\beta_0=-.956+\log(280/256)=-.866388\ldots$，主文记为 $-.8664$，医院 MDCEV 分配保持不变；评价层按 $G=280$ 重算。\\
结构参数 & $.15,1.20,.45,.40,.60$ 按“临床风险与 Hb 缺口为主、压力交互为辅”设置相对强度；$.08,.12$ 分别控制日级与患者扰动。\\
\end{tabularx}
\end{center}

\vfill
{\fontsize{7.8}{9.2}\selectfont\color{muted}
临床阈值参考：Carson et al. (2023), \textit{Red Blood Cell Transfusion: 2023 AABB International Guidelines}, \textit{JAMA} 330(19):1892--1902。本文在 7--8 g/dL 范围内以 $70+10r_i$ 连续参数化。}

\end{document}
